% META: {"id": "TNSI-ALGO-EVAL-B-corrige", "chapitre": "TNSI-ALGORITHMIQUE", "type_objet": "correction", "evaluation_ref": "TNSI-ALGO-EVAL-B", "capacites": ["T-ALGO-02A", "T-ALGO-02C", "T-ALGO-03"], "mode_creation": "ex_nihilo", "sources_inspiration": [], "status": "approved"}
% BEGIN-VERIFY
% def bfs(graphe, depart):
%     visites = [depart]
%     file = [depart]
%     while file:
%         sommet = file.pop(0)
%         for voisin in graphe[sommet]:
%             if voisin not in visites:
%                 visites.append(voisin)
%                 file.append(voisin)
%     return visites
% graphe = {
%     "P": ["Q", "R"], "Q": ["P", "S"], "R": ["P", "S"], "S": ["Q", "R"],
% }
% assert bfs(graphe, "P") == ["P", "Q", "R", "S"]
% def somme_dpr(liste):
%     if len(liste) == 1:
%         return liste[0]
%     milieu = len(liste) // 2
%     return somme_dpr(liste[:milieu]) + somme_dpr(liste[milieu:])
% assert somme_dpr([1, 2, 3, 4, 5]) == 15
% END-VERIFY
\begin{corrige}{TNSI-ALGO-EVAL-B-EX1}
\textbf{Q1.} Depuis \texttt{"P"} : on enfile \texttt{"P"}, on le défile et on enfile ses voisins \texttt{"Q"} puis \texttt{"R"} ; on défile \texttt{"Q"}, son voisin \texttt{"S"} n'est pas encore visité, on l'enfile ; on défile \texttt{"R"}, son voisin \texttt{"S"} est déjà visité. Ordre de visite : \texttt{P, Q, R, S}.

\textbf{Q2.} Oui : \texttt{P} $\to$ \texttt{Q} $\to$ \texttt{S} $\to$ \texttt{R} $\to$ \texttt{P} est un cycle (chaque arête utilisée existe bien dans le graphe, et on revient au sommet de départ sans repasser par une même arête).
\end{corrige}

\begin{corrige}{TNSI-ALGO-EVAL-B-EX2}
\textbf{Q1.} Diviser (découper le problème en sous-instances plus petites de même nature), régner (résoudre récursivement chaque sous-instance), combiner (assembler les solutions partielles en la solution globale).

\textbf{Q2.}
\begin{python}
def somme_dpr(liste):
    if len(liste) == 1:
        return liste[0]
    milieu = len(liste) // 2
    return somme_dpr(liste[:milieu]) + somme_dpr(liste[milieu:])
\end{python}
\end{corrige}
